\chapter{Einleitung}
\label{cha:Einleitung}

%\textit{(Gemeinsames Kapitel: Motivation, Projektziel der Phase 2, Übersicht über das Gesamtsystem ESP32 + Raspi 5 + BTS7960.)}

%\section{Motivation und Zielsetzung}

%Technische Demonstratoren stellen ein wichtiges Instrument dar, um komplexe Technologien anschaulich zu vermitteln und das Interesse potenzieller Studierender für technische Studiengänge zu wecken. Insbesondere auf Hochschulmessen oder Informationsveranstaltungen bieten interaktive Systeme die Möglichkeit, Besucher bereits aus größerer Entfernung auf einen Messestand aufmerksam zu machen und gleichzeitig einen niedrigschwelligen Einstieg in Gespräche mit Studierenden und Mitarbeitenden zu schaffen.
%
%Im Rahmen einer vorherigen Studienarbeit wurde hierfür ein erster funktionsfähiger Demonstrator entwickelt. Obwohl die grundlegenden Fahr- und Steuerungsfunktionen erfolgreich umgesetzt werden konnten, bestanden insbesondere hinsichtlich des Designs, der Mensch-Maschine-Interaktion sowie der technischen Zuverlässigkeit noch erhebliche Verbesserungspotenziale.
%
%Ziel dieser Studienarbeit ist daher die konstruktive, funktionale und gestalterische Weiterentwicklung des bestehenden Roboters. Neben der Entwicklung eines modernen Erscheinungsbildes sollen zusätzliche Interaktionsmöglichkeiten geschaffen und die Einsatzfähigkeit des Roboters unter realen Messebedingungen verbessert werden. 



\section{Motivation und Zielsetzung}

Technische Demonstratoren stellen ein wichtiges Instrument dar, um komplexe Technologien anschaulich zu vermitteln und das Interesse potenzieller Studierender für technische Studiengänge zu wecken. Insbesondere auf Hochschulmessen oder Informationsveranstaltungen bieten interaktive Systeme die Möglichkeit, Besucher bereits aus größerer Entfernung auf einen Messestand aufmerksam zu machen und gleichzeitig einen niedrigschwelligen Einstieg in Gespräche mit Studierenden und Mitarbeitenden zu schaffen.

Im Rahmen einer vorherigen Studienarbeit wurde hierfür ein erster funktionsfähiger Demonstrationsroboter entwickelt. Obwohl die grundlegenden Fahr- und Steuerungsfunktionen erfolgreich umgesetzt werden konnten, bestanden insbesondere hinsichtlich des mechanischen Aufbaus, des äußeren Erscheinungsbildes, der Mensch-Roboter-Interaktion sowie der technischen Zuverlässigkeit noch erhebliche Verbesserungspotenziale.

Ziel dieser Studienarbeit ist die konstruktive, funktionale und gestalterische Weiterentwicklung des bestehenden Roboters. Hierzu zählen die Entwicklung eines modernen Gehäusedesigns mit humanoiden Gestaltungselementen, die Verbesserung der Standfestigkeit des Fahrwerks sowie die Integration beweglicher Roboterarme und einer Hinderniserkennung zur Erhöhung der Betriebssicherheit. Darüber hinaus soll die Mensch-Roboter-Interaktion durch ein sprachbasiertes Dialogsystem erweitert werden, um eine natürlichere Kommunikation mit den Besuchern zu ermöglichen. 
Abschließend werden die entwickelten Konzepte und die umgesetzten Erweiterungen im Rahmen zweier Messeeinsätze unter realen Einsatzbedingungen validiert und hinsichtlich ihrer Praxistauglichkeit bewertet.


%\section{Systemübersicht und Hardware-Architektur}
%Beschreibung des Differential Drive Roboters...


%
%\section{Stand zu Projektbeginn}
%
%Die vorliegende Studienarbeit baut auf einer zuvor durchgeführten Studienarbeit auf, in der ein erster funktionsfähiger Demonstrationsroboter entwickelt wurde.
%Im Mittelpunkt dieser ersten Projektphase standen die Entwicklung eines Gesamtkonzepts, die Auswahl geeigneter Hard- und Softwarekomponenten sowie die Validierung grundlegender Funktionen des Roboters. \cite{Hadutoglu2026}
%
%Im Rahmen der vorangegangenen Arbeit wurde zunächst eine Anforderungsanalyse durchgeführt, auf deren Grundlage verschiedene Konstruktions- und Antriebskonzepte bewertet wurden. Mithilfe einer Trade-off-Analyse fiel die Entscheidung auf einen dreirädrigen Roboter mit Differenzialantrieb, da dieses Konzept eine hohe Wendigkeit sowie einen vergleichsweise einfachen mechanischen Aufbau bietet. Gleichzeitig wurde eine modulare Systemarchitektur entwickelt, bei der die Bewegungssteuerung und die Benutzerinteraktion auf zwei getrennte Recheneinheiten aufgeteilt wurden. Während ein ESP32 die Motorsteuerung übernimmt, dient ein Raspberry Pi als Plattform für das Display und zukünftige Interaktionsfunktionen.  
%
%Neben der elektrischen und softwareseitigen Konzeption wurde bereits ein erster mechanischer Prototyp gefertigt. Dieser bestand überwiegend aus einer Holzkonstruktion und diente der Erprobung des Fahrwerks sowie der Validierung der ausgewählten Komponenten. Darüber hinaus wurde die Motorsteuerung vollständig implementiert, sodass der Roboter über eine Fernsteuerung bewegt werden konnte. Die entwickelte Differenzialsteuerung ermöglichte sowohl Geradeausfahrten als auch Drehungen auf der Stelle. Zusätzlich wurden Sicherheitsmechanismen, wie ein Failsafe bei Verbindungsabbruch, implementiert und erfolgreich getestet.  
%
%Ein weiterer Schwerpunkt der vorherigen Studienarbeit lag auf der Entwicklung der Mensch-Maschine-Interaktion. Hierfür wurde ein Display in das System integriert, auf dem animierte Augen sowie Statusinformationen dargestellt werden sollten. Die grundlegende Softwarearchitektur zur Ansteuerung des Displays sowie erste Animationen wurden erfolgreich umgesetzt und in Betrieb genommen. Die Integration weiterer Interaktionsmöglichkeiten sowie die endgültige Gestaltung des Roboters waren jedoch nicht Bestandteil der damaligen Arbeit.  
%
%Zum Abschluss der ersten Projektphase lag somit ein funktionsfähiger Demonstrator vor, dessen Fahrwerk, Antriebssystem, Steuerung sowie grundlegende Interaktionsfunktionen erfolgreich validiert werden konnten. Die mechanische Konstruktion entsprach jedoch noch einem funktionalen Holzprototyp und war weder für den dauerhaften Einsatz noch hinsichtlich Design, Gewicht oder Modularität optimiert. Ebenso fehlten die vollständige Integration des Kopfes, ein finales Gehäusedesign sowie weitere Sensorik und autonome Funktionen. Diese offenen Punkte bildeten den Ausgangszustand und zugleich die Grundlage für die im Rahmen dieser Studienarbeit durchgeführte Weiterentwicklung des Roboters.  


%\subsection{Ziele der Arbeit}
%\textcolor{red}{brauchen wir das wirklich?}
%Ausgehend vom Entwicklungsstand der vorherigen Studienarbeit verfolgt diese Arbeit folgende Ziele:
%
%\begin{itemize}
%	
%	\item Entwicklung eines modernen Gehäusedesigns mit humanoiden Gestaltungselementen,
%	
%	\item Erhöhung der Standfestigkeit des Fahrwerks,
%	
%	\item Entwicklung und Integration beweglicher Roboterarme,
%	
%	\item Integration einer Hinderniserkennung zur Erhöhung der Betriebssicherheit,
%	
%	\item Erweiterung der Mensch-Roboter-Interaktion durch ein sprachbasiertes Dialogsystem,
%	
%	\item Validierung des Gesamtsystems unter realen Einsatzbedingungen auf Hochschulmessen.
%	
%\end{itemize}

\section{Aufbau der Arbeit}

Nach der Einleitung werden zunächst die für das Verständnis dieser Studienarbeit erforderlichen theoretischen Grundlagen erläutert. Anschließend wird der Entwicklungsstand des Roboters zu Beginn der Arbeit dargestellt, aus dem die im weiteren Verlauf behandelten Weiterentwicklungen abgeleitet werden.

Die Entwicklung des Roboters gliedert sich in drei thematische Schwerpunkte. Das erste Kapitel (\ref{cha:Vorgehensweise}. \nameref{cha:Vorgehensweise}) behandelt die technische Weiterentwicklung des Roboters. Hier werden die mechanischen und elektronischen Konzepte entwickelt, anhand von Trade-off-Analysen bewertet und anschließend praktisch umgesetzt. Im zweiten Schwerpunkt stehen das Design sowie die Verkleidung des Roboters (Kapitel \ref{cha:Design}) im Vordergrund. Dabei werden die Entwicklung des äußeren Erscheinungsbildes, die Konstruktion der Bauteile sowie deren Fertigung beschrieben. Das dritte Kapitel widmet sich den \nameref{cha:Umsetzung und Ergebnisse} des Roboters. Hierzu zählen insbesondere die Entwicklung und Integration der Sprachinteraktion sowie weiterer Funktionen zur Mensch-Roboter-Interaktion.

In allen drei Themenbereichen wird zunächst die Konzeption der jeweiligen Lösungsansätze vorgestellt und begründet. Anschließend wird deren praktische Umsetzung sowie die Integration in das Gesamtsystem beschrieben.

Den Abschluss der Entwicklungsarbeiten bildet die Validierung des Roboters unter realen Einsatzbedingungen. Hierzu wird das Gesamtsystem im Rahmen zweier Messeeinsätze erprobt und hinsichtlich Funktionalität, Zuverlässigkeit und Benutzerakzeptanz bewertet. Die dabei gewonnenen Erkenntnisse dienen als Grundlage für die abschließende Zusammenfassung der Arbeit sowie für den Ausblick auf mögliche zukünftige Weiterentwicklungen.

\section{Projektorganisation}
%Das Projkékt wurde von vier Mitgliedern durchgeführt. Faik Hadutoglu, Simon Steiger, Hai Nam Nguyen und Angus Bach. Dabei übernahm jeder einige Kernaufgaben, dazwischen fand geeinsamme Zusammenarbeit statt. In Tabelle \ref*{tab:RASIC} ist die Rollenverteilung in einem RASIC gezeigt.  Die übergänge wurden in der Umsetztung jedoch nicht besonder

Die Weiterentwicklung des Roboters erfolgte im Rahmen eines vierköpfigen Projektteams. Das Team umfasst: Faik Hadutoglu, Simon Steiger, Hai Nam Nguyen und Angus Bach. Aufgrund der unterschiedlichen Schwerpunkte der Studienarbeit wurden die Aufgaben auf die einzelnen Teammitglieder verteilt. Gleichzeitig erfolgte die Entwicklung zahlreicher Baugruppen in enger Abstimmung, sodass insbesondere bei der Systemintegration eine enge Zusammenarbeit erforderlich war. 

Diese Arbeit Umfasst die Ergebnisse von Hai Nam Nguyen und Angus Bach. Die Ergebnisse von Faik Hadutoglu werden in einer seperaten Studienarbeit behandelt, in dieser Abrbeit liegt der Fokus auf des Softwareentwicklung für Fahr und Steuerfunktionen.

Die Rollenverteilung innerhalb des Projekts ist in Tabelle~\ref{tab:RASIC} anhand eines RASIC-Diagramms dargestellt. Das Projekt orientierte sich an agilen Projektstrukturierung. Es wurden zeiwochen Sprints genutzt, 

\begin{table}[htbp]
	\centering
	\begin{tabular}{lcccc}
	
	\toprule
	
	\textbf{Aufgabe} & \textbf{Faik} & \textbf{Simon} & \textbf{Hai Nam} & \textbf{Angus} \\
	
	\midrule
	
	Projektkoordination        & S & S & S & R \\
	
	Mechanische Konstruktion   & I & I & S & R \\
	
	Elektronik                 & S & S & I & R \\
	
	Softwareentwicklung        & R & S & I & S \\
	
	KI-Modell                 & S & R & I & I \\
	
	Design                 & I & I & R & S \\
	
%	Systemintegration          & R & R & R & A \\
	
%	Test und Inbetriebnahme    & R & R & R & A \\
	
	Dokumentation              & C & R & R & R \\
	
	\bottomrule
	
\end{tabular}

\begin{tabular}{ll}
	
%	\toprule
	
	\textbf{Kürzel} & \textbf{Bedeutung} \\
	
	\midrule
	
	R & Responsible (ausführend) \\
	
	A & Accountable (verantwortlich) \\
	
	S & Supporting (unterstützend) \\
	
	I & Informed (informiert) \\
	
	C & Consulted (konsultiert) \\
	\bottomrule
	
\end{tabular}
	\caption{RASIC zur Rollenverteilung}
	\label{tab:RASIC}
\end{table}

